\documentclass[11pt]{article}
\usepackage[margin=1.1in]{geometry}
\usepackage{times}
\usepackage[T1]{fontenc}
\usepackage{setspace}
\onehalfspacing
\usepackage{parskip}
\pagenumbering{gobble}

\title{You're All Caught Up}
\author{Flyxion}
\date{September 2026}

\begin{document}
\maketitle

For half a second, the feed tells the truth. There is nothing left. You have reached the actual end of the actual content that actually existed at the moment you asked, a small honest fact rendered in a friendly sans-serif font, accompanied sometimes by an illustration of a cheerful animal resting, as if to say the day's obligations have been met and rest has been earned.

Then more appears. Not new content that arrived in the interval, though the interface is careful never to specify that it wasn't, but content that was always available and simply hadn't been shown yet, held back so the caught-up moment could be sold to you honestly the first time and then quietly withdrawn the second. The satisfaction was real for exactly as long as it took the system to decide it had been expressed for a plausible enough duration, and no longer, because a feed that let the caught-up feeling last would eventually have to confront a user who closed the app and did not reopen it, which is the one outcome the architecture exists to prevent.

Completion, in most of human experience, is a stable state. A meal finishes. A book ends. A task, once done, stays done until a new one is chosen. Here completion is rendered as a loading state, a screen you pass through rather than arrive at, indistinguishable in duration from the buffering icon it resembles, because the two are doing the same job: filling the gap between one piece of content and the next in a way that keeps you present for both. You were told, correctly, that you were caught up. You were not told that being caught up was never going to be allowed to be the end of anything.

\end{document}
